\subsection{LINGUAGEM DE PROGRAMAÇÃO: JAVA}

\en{Java} é uma linguagem de programação orientada a objetos, presente em diversos espaços. Os autores do \en{Java} organizaram os objetivos da linguagem em volta dos princípios: simplicidade, orientação a objetos, distribuição, robustez, segurança, neutralidade arquitetural, portabilidade, linguagem interpretada, performance, \en{multithreaded} e dinâmicidade. Como linguagem de programação, \en{Java} é frequentemente comparado com \en{C++}, onde ele se destaca na comparação por ter: \en{garbage collector} (não é necessário excluir as variáveis manualmente), o gerenciamento de apontadores de endereços de memória automático (no \en{C++} é algo que precisa ser levado em consideração no processo de construção código) e sua sintaxe mais amigável \cite{corejavafundamentals}.

No ambiente de sistemas distribuídos com \en{Java} pode-se destacar, como exemplo, o \en{CORBA} (\en{Common Object Request Broker Architecture}), esta tecnologia possibilita que sistemas totalmente distintos consigam operar em conjunto; \en{CORBA} complementa o \en{Java} por oferecer um ambiente de objetos distribuídos, a comunicação entre eles ocorre através de interfaces de programação \cite{oraclecorbajava}. Desse modo, é notável a versatilidade do \en{Java} como linguagem de programação; isso aliado à sua compatibilidade com o \en{Kotlin}, faz com que a escolha seja precisa para o contexto de desenvolvimento dos \en{middlewares} visados por esse texto. 
